\subsection{The prescribed bounded open-boundary window}
\label{sec:quadratic_prescribed_obc_window}

\begin{lemma}[Physical prefix bases of the prescribed minimal representation]
    \label{lem:quadratic_prescribed_prefix_basis}
    \lean{OBCChainTensor.linearIndependent_physicalPrefixMatrix,
        OBCChainTensor.cutColumnSpace_eq_span_physicalPrefixMatrix,
        OBCChainTensor.physicalPrefixBasis,
        OBCChainTensor.physicalPrefixBasis_apply,
        OBCChainTensor.toMatrix_cutSlice_physicalPrefixBasis}
    \leanok
    At a minimal bond of any prescribed OBC representation, its own
    physical prefix columns form a basis of the actual physical cut
    column space. Between two minimal bonds, the matrix of physical
    coordinate slicing in these bases is the original supplied site
    matrix.
\end{lemma}
\begin{proof}
    \leanok
    The actual coefficient matrix factors through the prescribed
    prefix and suffix factors. Its rank equals the bond dimension,
    so the prefix columns are independent. The physical column
    space is contained in their span, and equality of dimensions
    gives equality of spaces. The original adjacent-prefix recursion
    then identifies each supplied site matrix with physical slicing.
\end{proof}

\begin{lemma}[Derived Perron normalization inside the same physical window]
    \label{lem:quadratic_perron_normalization}
    \lean{MPSTensor.exists_unital_gauge_of_isNBlkInjective,
        MPSTensor.exists_apply_ne_zero_of_mpv_ne_zero,
        MPSTensor.quadraticWindow_cutCoefficientRank_eq_of_isNBlkInjective,
        MPSTensor.exists_quadraticSolution_physicalCutWindow_of_isNBlkInjective}
    \leanok
    \uses{lem:quadratic_virtual_cut_slicing,
        lem:quadratic_actual_physical_window_solution}
    Let $A$ satisfy C1 at length $L_0$ and have a nonzero letter.
    There are an invertible $X$ and a scalar $\zeta\ne0$ such that
    $B_i=\zeta X A_iX^{-1}$ is row normalized. C1 holds for the
    original tensor at every length at least $L_0$. Every actual
    physical cut-basis window therefore admits a normalized quadratic
    solution, retaining its original physical coefficient spaces and
    local matrices. A nonzero actual vector of positive length
    supplies the required nonzero letter.
\end{lemma}
\begin{proof}
    \leanok
    C1 gives full vector spreading, hence irreducibility, including
    at zero C1 length. Perron normalization gives $B$, $X$ and
    $\zeta$. Transport C1 to $B$, extend it by row normalization,
    and transport it back to $A$. If $Y_j,Z_j$ are the derived
    physical basis comparison matrices, use
    \begin{align}
        \widehat Y_j=\zeta^{-j}(X\otimes I)Y_j,\qquad
        \widehat Z_j=\zeta^{j+1}Z_j(X^{-1}\otimes I).
    \end{align}
    The scalar factor at a site is $\zeta$ and the amplified
    conjugation gives $B_i\otimes I$; at an internal inverse equation
    both cancel. No canonical form of $A$ is assumed.
\end{proof}

\begin{lemma}[A quadratic solution over the prescribed OBC window]
    \label{lem:quadratic_prescribed_obc_window_solution}
    \lean{MPSTensor.prescribedOBC_window_bondDim_eq,
        MPSTensor.prescribedOBCWindowBondEquiv,
        MPSTensor.quadraticPrescribedOBCWindow,
        MPSTensor.prescribedOBCWindowCutBases,
        MPSTensor.quadraticPrescribedOBCWindow_eq_physicalCutWindow,
        MPSTensor.exists_quadraticSolution_prescribedOBCWindow,
        MPSTensor.exists_quadraticSolution_prescribedOBCWindow_of_mpv_ne_zero}
    \leanok
    \uses{lem:quadratic_prescribed_prefix_basis,
        lem:quadratic_injective_physical_cut,
        lem:quadratic_perron_normalization}
    Let $A$ be a TI tensor satisfying C1 at length $L_0$, with
    $N>2L_0+D^4$ and a nonzero actual length-$N$ vector. Let $B$ be any
    prescribed OBC representation of its actual length-$N$ state,
    with all bonds bounded by $D^2$. Then every bond at the
    $D^4+1$ cuts of the reconstruction window has dimension exactly
    $D^2$. The original site matrices of $B$ in this window, expressed
    in the resulting doubled coordinates, admit a solution of the
    quadratic system. No solution, comparison gauge, or bond equality
    is supplied. The prescribed representation may in particular
    be the canonical one; uniqueness is unnecessary.
\end{lemma}
\begin{proof}
    \leanok
    At each window cut, C1 forces physical rank $D^2$. This rank is
    at most the original bond dimension, which is at most $D^2$.
    Hence the bond is minimal. The preceding lemma derives physical
    cut bases from the original prefix columns and identifies the
    original sites with physical slicing. The solution construction
    for any physical cut-basis window, with internally derived
    Perron normalization, therefore applies to this very
    prescribed representation. Only the forced coordinate
    identification is made; the representation is not replaced.
\end{proof}

\begin{lemma}[The original canonical OBC bonds are physical cut ranks]
    \label{lem:quadratic_canonical_obc_physical_ranks}
    \lean{OBCChainTensor.physicalPrefixMatrix_gram_step,
        OBCChainTensor.physicalPrefixMatrix_zero_gram,
        OBCChainTensor.physicalPrefixMatrix_gram_eq_density,
        OBCChainTensor.physicalSuffixMatrix_step,
        OBCChainTensor.physicalSuffixMatrix_last,
        OBCChainTensor.physicalSuffixMatrix_gram_step,
        OBCChainTensor.physicalSuffixMatrix_coisometry,
        OBCChainTensor.bondDim_eq_cutCoefficientRank_of_canonical,
        OBCChainTensor.coeff_ne_zero_of_faithful_density,
        OBCChainTensor.withBondBound,
        OBCChainTensor.withBondBound_tensor,
        OBCChainTensor.withBondBound_coeff}
    \leanok
    \uses{lem:quadratic_prescribed_prefix_basis}
    Let $B$ be the prescribed OBC tensor, with arbitrary common bond
    bound. Assume its source canonical equations
    \begin{align}
        \sum_i B_i^{[j]}B_i^{[j]\dagger}=I,\qquad
        \sum_i B_i^{[j]\dagger}\Lambda_j B_i^{[j]}=\Lambda_{j+1},
        \qquad\Lambda_0=I,\quad\Lambda_j>0.
    \end{align}
    At every cut, its actual prefix Gram is $\Lambda_j$ and its
    actual suffix Gram is the identity. Its original bond dimension
    is therefore exactly the physical cut rank. The represented
    vector is nonzero. Replacing only a verified common upper bound
    retains every original bond space, site matrix and coefficient.
\end{lemma}
\begin{proof}
    \leanok
    The physical extension bijections give the forward prefix and
    backward suffix Gram recursions. Induct from the two scalar
    boundaries. Faithfulness makes the prefix columns independent.
    If the physical coefficient matrix is $M=FG$, the suffix identity
    $GG^\dagger=I$ gives $MG^\dagger=F$, hence equality of the physical
    rank and the original bond dimension. Faithfulness at the last
    scalar bond also excludes a zero represented vector.
\end{proof}

\begin{lemma}[Quadratic existence for the original canonical OBC window]
    \label{lem:quadratic_original_canonical_obc_window}
    \lean{MPSTensor.canonicalOBC_bondDim_le_of_mpv_eq,
        MPSTensor.canonicalOBCWithSquaredBound,
        MPSTensor.quadraticCanonicalOBCWindow,
        MPSTensor.quadraticCanonicalOBCWindow_apply,
        MPSTensor.exists_quadraticSolution_canonicalOBCWindow}
    \leanok
    \uses{lem:quadratic_canonical_obc_physical_ranks,
        lem:quadratic_prescribed_obc_window_solution}
    Let the original canonical OBC tensor represent the actual
    length-$N$ state of a TI tensor $A$ satisfying C1 at $L_0$, with
    $N>2L_0+D^4$. Its original window admits a solution of the
    quadratic system. No common bond bound by $D^2$, nonzero-state
    premise, canonical TI representative, comparison gauge, window
    identity or quadratic solution is supplied. OBC uniqueness is
    unnecessary; the source canonical equations above suffice.
\end{lemma}
\begin{proof}
    \leanok
    The preceding lemma identifies the original bond dimensions
    with physical cut ranks. Opening a periodic trace contraction
    gives the bound $D^2$ for every such rank. This changes only the
    common upper bound recorded for the original rectangular tensors.
    The state is nonzero by faithfulness. Apply the prescribed-window
    theorem, with its internally derived Perron normalization and
    window gauges. The resulting window entries are the original
    supplied site entries in the proved doubled coordinates.
\end{proof}

\subsection{Endpoint gauges of an arbitrary solution}
\label{sec:quadratic_endpoint_gauges}

\begin{lemma}[A spanning length inside the window]
    \label{lem:quadratic_window_spanning_length}
    \lean{MPSTensor.finrank_wordSpan_le_one_of_krausRank_le_one,
        MPSTensor.exists_isNBlkInjective_le_of_canonical}
    \leanok
    \uses{def:pgvwc07_ti_canonical_representation, def:l_blk_injective}
    Let $A$ be a canonical tensor of bond dimension $D\ge2$ satisfying
    condition C1. Then the words of some length $\ell$ with
    $1\le\ell\le D^4-2$ in $A$ span $\MN{D}$.
\end{lemma}
\begin{proof}
    \leanok
    Condition C1 leaves one block of weight $\lambda>0$, and the unital
    block $\lambda^{-1}A$ has the same word spans as $A$. Its adjoint
    family is trace preserving and satisfies condition C1, hence is
    primitive. If its letters spanned a space of dimension at most one,
    every word of length $n$ would be a multiple of the $n$th power of a
    single matrix, and no word span could be the $D^2$-dimensional
    algebra $\MN{D}$. So the letters span a space of dimension $d'\ge2$,
    and the quantum Wielandt inequality bounds the least spanning length
    by $(D^2-d'+1)D^2\le D^4-D^2\le D^4-2$. Conjugate transposition
    reverses words, so the same length serves for $A$.
\end{proof}

\begin{lemma}[Invertibility of the left endpoint gauge]
    \label{lem:quadratic_gauged_endpoint}
    \lean{MPSTensor.evalWord_cons_kronecker_eq_of_gauged,
        MPSTensor.isUnit_gauge_zero_of_isNBlkInjective}
    \leanok
    \uses{def:pgvwc07_quadratic_system}
    Let $A,B\in\MN{D}^d$ and let $G_j,H_j$ be matrices on the doubled bond
    space with
    \begin{align}
        B_i\otimes\Id & =G_j(A_i\otimes\Id)H_j \qquad(0\le j\le m-1),
        \label{eq:quadratic_gauged_endpoint_system}                   \\
        G_{j+1}H_j    & =\Id \qquad(0\le j\le m-2).
        \label{eq:quadratic_gauged_endpoint_inverse}
    \end{align}
    A nonempty word $w$ placed at the sites $a,\ldots,a+|w|-1$ of the
    window satisfies $B_w\otimes\Id=G_a(A_w\otimes\Id)H_{a+|w|-1}$. If the
    words of some length $1\le\ell\le m-2$ in $A$ span $\MN{D}$, then $G_0$
    is invertible.
\end{lemma}
\begin{proof}
    \leanok
    The word identity follows by induction on the word, cancelling
    $H_aG_{a+1}=\Id$ between consecutive letters. Choose $c_w$ with
    $\sum_wc_wA_w=\Id$ over the words of length $\ell$ and put
    $X=\sum_wc_wB_w$. Placing these words at the sites $1,\ldots,\ell$ and
    at the sites $0,\ldots,\ell-1$ gives
    $X\otimes\Id=G_1H_\ell=G_0H_{\ell-1}$. Both $G_1$ and $H_\ell$ have
    one-sided inverses by (\ref{eq:quadratic_gauged_endpoint_inverse}),
    since $\ell+1\le m-1$, so $G_0H_{\ell-1}$ is invertible and so is $G_0$.
\end{proof}

\begin{lemma}[The complete chain for an arbitrary solution]
    \label{lem:quadratic_gauged_complete_chain}
    \lean{MPSTensor.exists_intertwiner_of_gauged,
        MPSTensor.exists_intertwiner_of_physicalCutWindow}
    \leanok
    \uses{lem:quadratic_gauged_endpoint, lem:quadratic_window_spanning_length,
        lem:pgvwc07_complete_chain_intertwiner,
        lem:quadratic_actual_physical_window_solution}
    In the setting of Lemma~\ref{lem:quadratic_gauged_endpoint} with
    $m=D^4$, let $G_0$ be invertible and let some $B_i$ be nonzero. Then
    there are a nonzero $R\in\MN{D}$ and a scalar $x\ne0$ with
    $RB_i=\tfrac1xA_iR$ for every $i$. In particular, let $A$ be a canonical
    tensor satisfying condition C1 at $L_0$, let $N>2L_0+D^4$, and let the
    window consist of the physical slicing matrices in any bases of the
    actual cut spaces of the length-$N$ state. Then every solution of the
    quadratic system over this window has such an intertwiner between its
    tensor and $A$.
\end{lemma}
\begin{proof}
    \leanok
    Each $H_j$ is nonzero, since otherwise every $B_i$ would vanish. The
    matrices $G_0^{-1},H_0,\ldots,H_{D^4-1}$ at the $D^4+1$ cuts satisfy
    the chain relations, the first from
    (\ref{eq:quadratic_gauged_endpoint_system}) at $j=0$ and the others from
    $H_jG_{j+1}=\Id$, so Lemma~\ref{lem:pgvwc07_complete_chain_intertwiner}
    applies. For the physical window, write the window matrices as
    $P_j^{-1}(A_i\otimes\Id)P_{j+1}$ by
    Lemma~\ref{lem:quadratic_actual_physical_window_solution} and put
    $G_j=Y_jP_j^{-1}$ and $H_j=P_{j+1}Z_j$ for a solution $(B,Y_j,Z_j)$.
    The normalization equation makes some $B_i$ nonzero. If $D\ge2$,
    Lemmas~\ref{lem:quadratic_window_spanning_length}
    and~\ref{lem:quadratic_gauged_endpoint} make $G_0$ invertible. If $D=1$,
    the gauge equation at the single site expresses a nonzero $B_i$ through
    the scalar $G_0$, which is therefore invertible.
\end{proof}
